\subsection{Trazabilidad de los controles y condiciones de aceptación}\label{sec:sd8-trazabilidad-p08}
Cada riesgo se vincula con su requisito, decisión o supuesto de origen, el componente o proceso afectado, el paquete e hito de ejecución, su propietario, la acción de tratamiento y la prueba prevista. Los registros de riesgos, acciones y escenarios permiten seguir esas relaciones hasta T-16. La evidencia de aceptación se localiza en SD9/T-13/T-17 y la asignación de recursos en SD7/T-15/T-18. Los protocolos PC-01 a PC-09 desarrollan las pruebas de continuidad; cada acción restante requiere su verificación funcional específica.

T-12 identifica la sección de SD8 o ficha de T-16 que desarrolla la respuesta, junto con el componente y la prueba prevista. Las remisiones se cotejan con la versión vigente de los documentos. La aceptación de un requisito exige evidencia de su diseño y funcionamiento; su identificación como riesgo no sustituye esa evidencia. \parencites[numeral~1.5, p.~4]{sd8-bt}[arts.~46--47, p.~28]{sd8-ba}

Antes de autorizar la puesta en producción, los responsables deben acreditar autonomía e interacción de los terminales, permisos durante aislamiento, habilitación y capacidad de recuperación ante desastre, restauración por dominio, suplencia del legado y disponibilidad de recursos. Calidad verifica los resultados y la Contraparte Técnica formaliza la recepción que corresponda. Los cálculos, reservas y protocolos de esta oferta establecen el tratamiento propuesto; los resultados de las pruebas se incorporarán al expediente de aceptación durante la ejecución.
